%Daten zur Institution

%\input{Dozentdaten}

%\renewcommand{\fachbereich}{Fachbereich}

%\renewcommand{\dozent}{Prof. Dr. . }

%Klausurdaten

\renewcommand{\klausurgebiet}{ }

\renewcommand{\klausurtyp}{ }

\renewcommand{\klausurdatum}{ . 20}

\klausurvorspann {\fachbereich} {\klausurdatum} {\dozent} {\klausurgebiet} {\klausurtyp}

%Daten für folgende Punktetabelle

\renewcommand{\aeins}{  3  }

\renewcommand{\azwei}{  3   }

\renewcommand{\adrei}{  8  }

\renewcommand{\avier}{  7  }

\renewcommand{\afuenf}{  8  }

\renewcommand{\asechs}{  3  }

\renewcommand{\asieben}{  5  }

\renewcommand{\aacht}{  5  }

\renewcommand{\aneun}{  6  }

\renewcommand{\azehn}{  4  }

\renewcommand{\aelf}{  12  }

\renewcommand{\azwoelf}{  64  }

\renewcommand{\adreizehn}{     }

\renewcommand{\avierzehn}{    }

\renewcommand{\afuenfzehn}{    }

\renewcommand{\asechzehn}{    }

\renewcommand{\asiebzehn}{    }

\renewcommand{\aachtzehn}{    }

\renewcommand{\aneunzehn}{    }

\renewcommand{\azwanzig}{    }

\renewcommand{\aeinundzwanzig}{    }

\renewcommand{\azweiundzwanzig}{    }

\renewcommand{\adreiundzwanzig}{    }

\renewcommand{\avierundzwanzig}{    }

\renewcommand{\afuenfundzwanzig}{    }

\renewcommand{\asechsundzwanzig}{    }

\punktetabelleelf

\klausurnote

\newpage

\setcounter{section}{0}

__NOEDITSECTION__

\inputaufgabeklausurloesung
{3}

{

Definiere die folgenden
\zusatzklammer {kursiv gedruckten} {} {} Begriffe.
\aufzaehlungsechs{Ein
\stichwort {Normalenfeld} {}
$N$ auf einer differenzierbaren Hyperfläche

\mavergleichskette
{\vergleichskette
{ Y
}
{ \subseteq }{ \R^n
}
{  }{
}
{    }{
}
{   }{
}
}
{}{}{.}

}{Ein
\stichwort {überdeckungskompakter} {}
topologischer Raum.

}{Der
\stichwort {Kotangentialraum} {}
in einem Punkt

\mavergleichskette
{\vergleichskette
{  P
}
{ \in }{ M
}
{  }{}
{    }{}
{   }{}
}
{}{}{}
einer
\definitionsverweis {differenzierbaren Mannigfaltigkeit}{}{}
$M$.

}{Ein
\stichwort {regulärer Punkt} {}

\mathl{Q \in L}{} zu einer
\definitionsverweis {differenzierbaren Abbildung}{}{}
\maabbdisp {\varphi} { L } { M
} {}
zwischen
\definitionsverweis {differenzierbaren Mannigfaltigkeiten}{}{}
\mathkor {} {L} {und} {M} {.}

}{Das
\stichwort {Produkt} {}
von differenzierbaren Mannigfaltigkeiten
\mathkor {} {M} {und} {N} {.}

}{Ein
\stichwort {Zusammenhang} {}
auf einem differenzierbaren Vektorbündel
\maabb {} {E} {M
} {}
über einer differenzierbaren Mannigfaltigkeit $M$.
}

}
{

\aufzaehlungsechs{Ein
Normalenfeld
auf $Y$ ist ein auf einer offenen Umgebung

\mavergleichskette
{\vergleichskette
{ Y
}
{ \subseteq }{ U
}
{ \subseteq }{ \R^n
}
{    }{
}
{   }{
}
}
{}{}{}
definiertes
\definitionsverweis {stetiges}{}{}
\definitionsverweis {Vektorfeld}{}{}
\maabb {F} {U} {\R^n
} {}
mit

\mavergleichskettedisp
{\vergleichskette
{ F(P)
}
{ \in} { N_PY
}
{ } {
}
{ } {
}
{ } {
}
}
{}{}{}
für alle

\mavergleichskette
{\vergleichskette
{ P
}
{ \in }{ Y
}
{  }{
}
{    }{
}
{   }{
}
}
{}{}{.}
}{Ein
\definitionsverweis {topologischer Raum}{}{}
$X$ heißt überdeckungskompakt, wenn es zu jeder offenen Überdeckung

\mathdisp {X= \bigcup_{i \in I} U_i \, \, \, \text{ mit } U_i \text{ offen und einer beliebigen Indexmenge }} {  }

eine endliche Teilmenge
\mathl{J \subseteq I}{} derart gibt, dass

\mavergleichskettedisp
{\vergleichskette
{ X
}
{ =} {\bigcup_{i \in J} U_i
}
{ } {
}
{ } {
}
{ } {
}
}
{}{}{}
ist.
}{Unter dem Kotangentialraum an $P$ versteht man den
\definitionsverweis {Dualraum}{}{}
des
\definitionsverweis {Tangentialraumes}{}{}

\mathl{T_PM}{} an $P$.
}{Der Punkt

\mavergleichskette
{\vergleichskette
{  Q
}
{ \in }{ L
}
{  }{
}
{    }{
}
{   }{
}
}
{}{}{}
heißt regulär für $\varphi$, wenn die
\definitionsverweis {Tangentialabbildung}{}{}
\maabbdisp {T_Q(\varphi)} {T_QL} {T_{\varphi(Q)}M
} {}
im Punkt $Q$
\definitionsverweis {maximalen Rang}{}{}
besitzt.
}{Es seien
\mathkor {} {(U_i,U_i',\alpha_i, i \in I)} {und} {(V_j,V_j',\beta_j, j \in J)} {}
die
\definitionsverweis {Atlanten}{}{}
von
\mathkor {} {M} {und} {N} {.}
Dann nennt man den
\definitionsverweis {Produktraum}{}{}

\mathl{M \times N}{} versehen mit den
\definitionsverweis {Karten}{}{}
\maabbdisp {\alpha_i \times \beta_j} {U_i \times V_j } { U_i' \times V_j'
} {}
\zusatzklammer {mit \mathlk{(i,j) \in I \times J}{} und \mathlk{U_i' \times V_j' \subseteq \R^m \times \R^n}{}} {} {}
das Produkt der Mannigfaltigkeiten
\mathkor {} {M} {und} {N} {.}
}{Unter einem
Zusammenhang
auf $E$ versteht man eine direkte Summenzerlegung des
\definitionsverweis {Tangentialbündels}{}{}

\mavergleichskette
{\vergleichskette
{ TE
}
{ = }{ V \oplus H
}
{  }{
}
{    }{
}
{   }{
}
}
{}{}{}
in zwei Untervektorbündel
\mathkor {} {V} {und} {H} {,}
wobei

\mavergleichskette
{\vergleichskette
{ V
}
{ \subseteq }{ TE
}
{  }{
}
{    }{
}
{   }{
}
}
{}{}{}
das
\definitionsverweis {Vertikalbündel}{}{}
ist.
}

 }

__NOEDITSECTION__

\inputaufgabeklausurloesung
{3}

{

Formuliere die folgenden Sätze.
\aufzaehlungdrei{Der
\stichwort {Satz über die Beziehung zwischen Krümmung und Einheitsnormalenvektor einer bogenparametrisierten ebenen Kurve} {.}}{Das
\stichwort {Theorema egregium} {.}}{Der
\stichwort {Satz von Green} {}
für den Flächeninhalt.}

}
{

\aufzaehlungdrei{\faktsituation {Es sei
\maabb {\gamma} { I } {\R^2
} {}
eine zweifach
\definitionsverweis {stetig differenzierbare}{}{}
\definitionsverweis {bogenparametrisierte}{}{}
\definitionsverweis {Kurve}{}{.}}
\faktfolgerung {Dann ist die
\definitionsverweis {Krümmung}{}{}
von $\gamma$ in $t$ gleich

\mavergleichskettedisp
{\vergleichskette
{ \kappa(t)
}
{ =} { \left\langle  \gamma^{\prime \prime}(t)  ,  N(t )    \right\rangle
}
{ =} { \pm \Vert { \gamma^{\prime \prime}(t) } \Vert
}
{ } {
}
{ } {
}
}
{}{}{,}
wobei

\mavergleichskettedisp
{\vergleichskette
{ N( t)
}
{ =} { \begin{pmatrix}  -\gamma_2'(t) \\ \gamma_1'(t)   \end{pmatrix}
}
{ } {
}
{ } {
}
{ } {
}
}
{}{}{}
ein Einheitsnormalenvektor in $\gamma(t)$ ist.}
\faktzusatz {}
\faktzusatz {}}{\faktsituation {Es sei

\mavergleichskette
{\vergleichskette
{ Y
}
{ \subseteq }{ \R^3
}
{  }{
}
{    }{
}
{   }{
}
}
{}{}{}
eine
\definitionsverweis {orientierte Fläche}{}{}
und sei
\maabbdisp {\varphi} { V } { Y
} {,}

\mavergleichskette
{\vergleichskette
{ V
}
{ \subseteq }{ \R^2
}
{  }{
}
{    }{
}
{   }{
}
}
{}{}{,}
eine zweifach differenzierbare lokale Parametrisierung von $Y$ mit den Parametern $u,v$. Es sei
\mathl{\left(  g_{ij}  \right)}{} die
\definitionsverweis {erste Fundamentalmatrix}{}{}
auf $V$.}
\faktfolgerung {Dann gilt für die
\definitionsverweis {Gaußsche Krümmung}{}{}
unter Verwendung der
\definitionsverweis {Christoffelsymbole}{}{}
die Beziehung

\mavergleichskettedisphandlinks
{\vergleichskettedisphandlinks
{ K
}
{ =} { { \frac{   1  }{ g_{11}  } }  \left(  \partial_2 \Gamma_{11}^2 - \partial_1 \Gamma_{12}^2 + \Gamma^1_{11} \Gamma^2_{12}  + \Gamma_{11}^2 \Gamma_{22}^2 - \Gamma_{12}^1 \Gamma_{11}^2-  \left(  \Gamma_{12}^2  \right)^2  \right)
}
{ } {
}
{ } {
}
{ } {
}
}
{}{}{.}}
\faktzusatz {}
\faktzusatz {}}{Es sei
\mathl{M \subseteq \R^2}{} eine
\definitionsverweis {kompakte}{}{}
\definitionsverweis {Mannigfaltigkeit mit Rand}{}{}

\mathl{\delta M}{.} Dann ist der Flächeninhalt von $M$ gleich

\mavergleichskettedisp
{\vergleichskette
{  \lambda^2(M)
}
{ =} {
\int_{ \partial M  }  x dy
}
{ =} { -
\int_{ \partial M  }  y dx
}
{ } {
}
{ } {
}
}
{}{}{.}}

 }

__NOEDITSECTION__

\inputaufgabeklausurloesung
{8}

{

Es sei

\mavergleichskettedisp
{\vergleichskette
{ f(x,y)
}
{ =} { ax^2+by^2+cxy
}
{ } {
}
{ } {
}
{ } {
}
}
{}{}{.}
Bestimme die
\definitionsverweis {Hauptkrümmungen}{}{}
und die
\definitionsverweis {Hauptkrümmungsrichtungen}{}{}
des
\definitionsverweis {Graphen}{}{}
zu $f$ im Nullpunkt.

}
{

Wir verwenden
[[Differenzierbare Funktion/Graph/Hyperfläche/Weingartenabbildung/Fakt|Lemma 4.9 (Differentialgeometrie (Osnabrück 2023))]],
wobei im Nullpunkt der Vorfaktor gleich $1$ ist und somit die
\definitionsverweis {Weingartenabbildung}{}{}
direkt durch die
\definitionsverweis {Hesse-Matrix}{}{}
beschrieben wird. Die Hesse-Matrix der Funktion ist

\mathdisp {\begin{pmatrix}  2a  &   c   \\ c &  2b \end{pmatrix}} { , }

das
\definitionsverweis {charakteristische Polynom}{}{}
davon ist

\mavergleichskettealign
{\vergleichskettealign
{ (t-2a)(t-2b)-c^2
}
{ =} { t^2 -2(a+b)t +4ab -c^2
}
{ =} { \left(  t-(a+b)  \right)^2 - (a+b)^2 +4ab -c^2
}
{ =} { \left(  t-(a+b)  \right)^2  - a^2 -b^2 +2ab -c^2
}
{ =} { \left(  t-(a+b)  \right)^2  - (a-b)^2 -c^2
}
}
{}
{}{.}
Die Eigenwerte
\zusatzklammer {und dies sind die Hauptkrümmungen} {} {}
sind also

\mavergleichskettedisp
{\vergleichskette
{ \lambda_{1,2}
}
{ =} { \pm \sqrt{ (a-b)^2+c^2 } +a+b
}
{ } {
}
{ } {
}
{ } {
}
}
{}{}{.}
Bei

\mavergleichskette
{\vergleichskette
{ a
}
{ = }{ b
}
{  }{
}
{    }{
}
{   }{
}
}
{}{}{}
und

\mavergleichskette
{\vergleichskette
{ c
}
{ = }{ 0
}
{  }{
}
{    }{
}
{   }{
}
}
{}{}{}
ist die Weingartenabbildung Multiplikation mit $2a$ und jeder Vektor ist ein Eigenvektor, sei dies jetzt ausgeschlossen. Für

\mavergleichskettedisp
{\vergleichskette
{ \lambda_{1}
}
{ =} { \sqrt{ (a-b)^2+c^2 } +a+b
}
{ } {
}
{ } {
}
{ } {
}
}
{}{}{}
ist der Kern der Matrix

\mathdisp {\begin{pmatrix}  \sqrt{(a-b)^2 +c^2} + b-a  &   -c  \\  -c &  \sqrt{(a-b)^2 +c^2} + a-b \end{pmatrix}} {  }

zu bestimmen, dieser ist

\mathdisp {\R \begin{pmatrix}  \sqrt{(a-b)^2 +c^2} + a-b  \\c   \end{pmatrix}} { . }

Eine Hauptkrümmungsrichtung ist also
\mathl{\begin{pmatrix}  \sqrt{(a-b)^2 +c^2} + a-b  \\c   \end{pmatrix}}{} mit der Hauptkrümmung
\mathl{\sqrt{ (a-b)^2+c^2 } +a+b}{.}

Entsprechend ergibt sich die Hauptkrümmungsrichtung
\mathl{\begin{pmatrix}  -c \\ \sqrt{(a-b)^2 +c^2} + a-b    \end{pmatrix}}{} mit der Hauptkrümmung
\mathl{-\sqrt{ (a-b)^2+c^2 } +a+b}{.}

 }

__NOEDITSECTION__

\inputaufgabeklausurloesung
{7}

{

Beweise den Existenzsatz für den Paralleltransport auf einer Hyperfläche

\mavergleichskette
{\vergleichskette
{ Y
}
{ \subseteq }{ \R^n
}
{  }{
}
{    }{
}
{   }{
}
}
{}{}{.}

}
{

Wir suchen nach einer differenzierbaren Abbildung
\maabbdisp {F} { I } { \R^n
} {,}
das die Bedingungnen
\aufzaehlungdrei{
\mavergleichskettedisp
{\vergleichskette
{ F(t)
}
{ \in} { T_{\gamma(t)} Y
}
{ } {
}
{ } {
}
{ } {
}
}
{}{}{}
für alle

\mavergleichskette
{\vergleichskette
{ t
}
{ \in }{ I
}
{  }{
}
{    }{
}
{   }{
}
}
{}{}{,}
}{
\mavergleichskettedisp
{\vergleichskette
{ F'(t)
}
{ \in} { N_{\gamma(t)} Y
}
{ } {
}
{ } {
}
{ } {
}
}
{}{}{}
für alle

\mavergleichskette
{\vergleichskette
{ t
}
{ \in }{ I
}
{  }{
}
{    }{
}
{   }{
}
}
{}{}{,}
}{
\mavergleichskettedisp
{\vergleichskette
{ F (t_0)
}
{ =} { v
}
{ } {
}
{ } {
}
{ } {
}
}
{}{}{}
}
erfüllt. Die erste Bedingung bedeutet, dass $F(t)$ senkrecht auf dem Normaleneinheitsvektors $N (\gamma(t))$ steht, also

\mavergleichskettedisp
{\vergleichskette
{ \left\langle F( t) ,  N( \gamma(t))   \right\rangle
}
{ =} { 0
}
{ } {
}
{ } {
}
{ } {
}
}
{}{}{.}
Daher ist auch

\mavergleichskettedisp
{\vergleichskette
{ 0
}
{ =} { \left\langle F(t)  ,  N( \gamma(t))   \right\rangle'
}
{ =} { \left\langle F'(t)  ,  N( \gamma(t))   \right\rangle +  \left\langle F(t)  ,  (N( \gamma(t)))'   \right\rangle
}
{ } {
}
{ } {
}
}
{}{}{.}
Die zweite Bedingung, dass $F'(t)$ im Normalenraum liegt, bedeutet

\mavergleichskettedisp
{\vergleichskette
{ F'(t) - \left\langle F'( t) ,  N( \gamma(t))   \right\rangle N(\gamma(t) )
}
{ =} { 0
}
{ } {
}
{ } {
}
{ } {
}
}
{}{}{}
für alle

\mavergleichskette
{\vergleichskette
{ t
}
{ \in }{ I
}
{  }{
}
{    }{
}
{   }{
}
}
{}{}{,}
was wir mithilfe der vorstehenden Gleichung zu

\mavergleichskettedisp
{\vergleichskette
{ F'(t) + \left\langle F( t) ,  (N( \gamma(t)))'   \right\rangle N(\gamma(t) )
}
{ =} { 0
}
{ } {
}
{ } {
}
{ } {
}
}
{}{}{}
für alle

\mavergleichskette
{\vergleichskette
{ t
}
{ \in }{ I
}
{  }{
}
{    }{
}
{   }{
}
}
{}{}{}
bzw. zu

\mavergleichskettedisp
{\vergleichskette
{ F'(t)
}
{ =} { - \left\langle F( t) ,  (N( \gamma(t)))'   \right\rangle N(\gamma(t) )
}
{ } {
}
{ } {
}
{ } {
}
}
{}{}{}
für alle

\mavergleichskette
{\vergleichskette
{ t
}
{ \in }{ I
}
{  }{
}
{    }{
}
{   }{
}
}
{}{}{}
umformen können. Hier steht eine gewöhnliche lineare Differentialgleichung erster Ordnung für $F$, die zusammen mit der Anfangsbedingung

\mavergleichskette
{\vergleichskette
{ F(t_0)
}
{ = }{ v
}
{  }{
}
{    }{
}
{   }{
}
}
{}{}{}
nach
[[Picard Lindelöf/Lokal Lipschitz/Lokale Existenz und Eindeutigkeit/Fakt|Satz 56.2 (Analysis (Osnabrück 2021-2023))]]
eine eindeutige Lösung besitzt. Es kann also höchstens eine Lösung der Ausgangsgleichung geben. Wenn $F$ die eindeutige Lösung der Differentialgleichung ist, so liegt in der Tat eine Lösung der Ausgangsgleichung vor.

 }

__NOEDITSECTION__

\inputaufgabeklausurloesung
{8}

{

Zeige, dass eine
\definitionsverweis {differenzierbare Mannigfaltigkeit}{}{}
$M$ der Dimension

\mavergleichskette
{\vergleichskette
{ n
}
{ \geq }{ 1
}
{  }{
}
{    }{
}
{   }{
}
}
{}{}{}
unendlich viele
\definitionsverweis {Diffeomorphismen}{}{}
\maabbdisp {\varphi} { M } { M
} {}
besitzt.

}
{

Wir betrachten Funktionen der Bauart
\maabbeledisp {\psi_c} {[0,1]} {[0,1]
} { x } { x + \varphi_c (x)
} {,}
mit

\mavergleichskettedisp
{\vergleichskette
{  \varphi_c(x)
}
{ =} { \begin{cases}  0 \text{ für } x \leq { \frac{   1  }{  3 } } \, , \\  c \left(  x- { \frac{   1  }{  3 } }   \right)^2 \left(  x- { \frac{   2 }{  3 } }   \right)^2 \text{ für } { \frac{   1  }{  3 } } < x < { \frac{   2 }{  3 } }  \, ,  \\  0 \text{ für } x \geq { \frac{   2 }{  3 } }  \, .  \end{cases}
}
{ } {
}
{ } {
}
{ } {
}
}
{}{}{}
mit einem Parameter

\mavergleichskette
{\vergleichskette
{ c
}
{ \in }{  \R_{\geq 0}
}
{  }{
}
{    }{
}
{   }{
}
}
{}{}{.}
Dabei ist $\varphi_c$ stetig und auch stetig differenzierbar, da die Nullstellen des Polynoms doppelt sind. Somit ist auch $\psi_c$ stetig differenzierbar. Für $c$ hinreichend klein ist $\psi_c$ streng wachsend und definiert eine stetig differenzierbare bijektive Abbildung des Einheitsintervalls in sich, die im unteren und im oberen Drittel die Identität ist.

Mit $\psi_c$ definieren wir Diffeomorphismen des offenen Einheitsballes in sich, durch
\maabbeledisp {} { U \left(   0 , 1  \right) } { U \left(   0 , 1  \right)
} { \begin{pmatrix}  x_1  \\\vdots \\  x_n   \end{pmatrix} } { \psi_c \left(  \sqrt{  x_1^2 + \cdots + x_n^2 }  \right) \begin{pmatrix}  x_1  \\\vdots \\  x_n   \end{pmatrix}
} {.}
Es wird also der Punkt abhängig vom Radius gestreckt, wobei auf
\mathl{U \left(   0 , { \frac{   1  }{  3 } }   \right)}{} und auf
\mathl{U \left(   0 , 1  \right) \setminus  U \left(   0 ,   { \frac{   2 }{  3 } }   \right)}{} die Identität vorliegt.

Wenn nun $M$ eine Mannigfaltigkeit der Dimension $\geq 1$ ist, so wählen wir eine Karte
\maabbdisp {\alpha} { U } {V
} {}
mit

\mavergleichskette
{\vergleichskette
{ V
}
{ \subseteq }{  \R^n
}
{  }{
}
{    }{
}
{   }{
}
}
{}{}{,}
wobei wir
\zusatzklammer {durch verkleinern} {} {}
davon ausgehen können, dass $V$
\zusatzklammer {ein offener Ball und dann auch} {} {}
der offene Einheitsball ist. Die konstruierten Diffeomorphismen auf $V$ liefern Diffeomorphismen auf $U$. Da diese für den äußeren Ball
\zusatzklammer {ab Radius ${ \frac{   2 }{  3 } }$} {} {}
die Identität sind, kann man diese Diffeomorphismen auf $M$ durch die Identität auf $M \setminus \alpha^{-1} \left(  U \left(   0 , { \frac{   2 }{  3 } }   \right)  \right)$ diffeomorph ausdehnen.

 }

__NOEDITSECTION__

\inputaufgabeklausurloesung
{3}

{

Es sei $M$ eine
\definitionsverweis {differenzierbare Mannigfaltigkeit}{}{}
und
\maabbdisp {f} { M } {\R
} {}
eine
\definitionsverweis {stetig differenzierbare Funktion}{}{.}
Die Funktion habe im Punkt

\mavergleichskette
{\vergleichskette
{ P
}
{ \in }{ M
}
{  }{
}
{    }{
}
{   }{
}
}
{}{}{}
ein
\definitionsverweis {lokales Extremum}{}{.}
Zeige

\mavergleichskettedisp
{\vergleichskette
{ (df)_P
}
{ =} { 0
}
{ } {
}
{ } {
}
{ } {
}
}
{}{}{.}

}
{

Es sei

\mavergleichskette
{\vergleichskette
{ P
}
{ \in }{  U
}
{ \subseteq }{ M
}
{    }{
}
{   }{
}
}
{}{}{}
ein
\definitionsverweis {Kartengebiet}{}{}
mit der Karte
\maabb {\alpha} { U } {V \subseteq \R^n
} {.}
Dann besitzt auch
\mathl{f \circ \alpha^{-1}}{} im Punkt $\alpha(P)$ ein lokales Extremum.
Daher ist nach
[[Lokales Extremum/Richtungsableitung/Totales Differential/Fakt|Fakt]] [[Kurs:Differentialgeometrie/Lokales Extremum/Richtungsableitung/Totales Differential/Fakt/Faktreferenznummer|*****]]  (2)
das
\definitionsverweis {totale Differential}{}{}
\maabbdisp {\left(D f \circ \alpha^{-1} \right)_{\alpha(P) }} {\R^n } {\R
} {}
die Nullabbildung. Daher ist

\mavergleichskette
{\vergleichskette
{ (df)_P
}
{ = }{ 0
}
{  }{
}
{    }{
}
{   }{
}
}
{}{}{.}

 }

__NOEDITSECTION__

\inputaufgabeklausurloesung
{5}

{

Wir betrachten die
$1$-\definitionsverweis {Differentialform}{}{}

\mavergleichskettedisp
{\vergleichskette
{ \omega
}
{ =} { -vdu+udv
}
{ } {
}
{ } {
}
{ } {
}
}
{}{}{}
auf der
$1$-\definitionsverweis {Sphäre}{}{}

\mavergleichskette
{\vergleichskette
{ S^1
}
{ \subset }{ \R^2
}
{  }{
}
{    }{
}
{   }{
}
}
{}{}{,}
wobei die Koordinaten des umgebenden $\R^2$ mit
\mathkor {} {u} {und} {v} {}
bezeichnet seien. Bestimme
\mathl{\pi^* \omega}{} unter der
\definitionsverweis {stetig differenzierbaren}{}{}
Abbildung
\maabbeledisp {\pi} {\R^2 \setminus \{0\}} { S^{1}
} {(x , y) } { { \frac{   1  }{  \sqrt{x^2 + y^2}  } } (x , y)
} {.}

}
{

Es ist

\mavergleichskettealign
{\vergleichskettealign
{ \pi^*du
}
{ =} { d (u \circ \pi)
}
{ =} { d   \left(  { \frac{   x  }{  \sqrt{x^2 + y^2}  } }  \right)
}
{ =} { { \frac{   \partial   \left(  { \frac{   x  }{  \sqrt{x^2 + y^2}  } }  \right)    }{  \partial   x  } }   dx + { \frac{   \partial   \left(  { \frac{   x  }{  \sqrt{x^2 + y^2}  } }  \right)    }{  \partial   y  } } dy
}
{ =} { { \frac{   \sqrt{x^2+y^2}     -x^2 (x^2 + y^2)^{-1/2}  }{  x^2+y^2  } }  dx   - { \frac{   xy }{  \sqrt{x^2+y^2}^3  } }  dy
}
}
{
\vergleichskettefortsetzungalign
{ =} { { \frac{   x^2+y^2 -x^2    }{  \sqrt{ x^2+y^2}^3  } }  dx  - { \frac{   xy }{  \sqrt{x^2+y^2}^3  } }  dy
}
{ =} { { \frac{   y^2  }{  \sqrt{ x^2+y^2}^3  } }  dx  - { \frac{   xy }{  \sqrt{x^2+y^2}^3  } }  dy
}
{ } {}
{ } {}
}
{}{}
und

\mavergleichskettealign
{\vergleichskettealign
{ \pi^*dv
}
{ =} { d (v \circ \pi)
}
{ =} { d \left(  { \frac{   y  }{  \sqrt{x^2 + y^2}  } }  \right)
}
{ =} { { \frac{   \partial   \left(  { \frac{   y  }{  \sqrt{x^2 + y^2}  } }  \right)   }{  \partial   x  } } dx + { \frac{   \partial   \left(  { \frac{   y  }{  \sqrt{x^2 + y^2}  } }  \right)    }{  \partial   y  } } dy
}
{ =} { { \frac{   -xy }{  \sqrt{x^2+y^2}^3  } } dx + { \frac{   \sqrt{x^2+y^2} -y^2 (x^2 + y^2)^{-1/2}  }{  x^2+y^2  } } dy
}
}
{
\vergleichskettefortsetzungalign
{ =} { { \frac{   -xy }{  \sqrt{x^2+y^2}^3  } } dx + { \frac{   x^2  }{  \sqrt{x^2+y^2}^3  } }  dy
}
{ } {
}
{ } {}
{ } {}
}
{}{.}
Somit ist

\mavergleichskettealigndrucklinks
{\vergleichskettealigndrucklinks
{ \pi^*(-vdu+udv)
}
{ =} { - (v \circ \pi)  \pi^*du + (u \circ \pi) \pi^*dv
}
{ =} { - { \frac{   y  }{  \sqrt{x^2+y^2}  } } \left(  { \frac{   y^2  }{  \sqrt{ x^2+y^2}^3  } }  dx  - { \frac{   xy }{  \sqrt{x^2+y^2}^3  } }  dy   \right) + { \frac{   x  }{  \sqrt{x^2+y^2}  } }  \left(  { \frac{   -xy }{  \sqrt{x^2+y^2}^3  } } dx + { \frac{   x^2  }{  \sqrt{x^2+y^2}^3  } } dy  \right)
}
{ =} { { \frac{   1 }{  \left(  x^2+y^2  \right)^2 } }  \left(  \left(  -y^3-x^2y  \right) dx + \left(  xy^2+x^3  \right) dy \right)
}
{ =} { { \frac{   1 }{  x^2+y^2  } }  \left(  -  y dx + x dy \right)
}
}
{}{}{.}

 }

__NOEDITSECTION__

\inputaufgabeklausurloesung
{5}

{

Beweise den Satz über die Berechnung des kanonischen Volumens auf einer riemannschen Mannigfaltigkeit $M$.

}
{

Gemäß der
[[Mannigfaltigkeit/Abzählbar/Positive Volumenform/Zugehöriges Maß/Definition|Definition  .]]
müssen wir die Differentialform
\mathl{\left(  \alpha^{-1}  \right)^*\omega}{} für jeden Punkt

\mavergleichskette
{\vergleichskette
{ Q
}
{ \in }{ V
}
{  }{
}
{    }{
}
{   }{
}
}
{}{}{}
berechnen. Diese Form besitzt die Gestalt
\mathl{c_Q dx_1 \wedge \ldots \wedge dx_n}{} und ist durch ihren Wert auf
\mathl{e_1 \wedge \ldots \wedge e_n}{} festgelegt. Es ist

\mavergleichskettedisphandlinks
{\vergleichskettedisphandlinks
{ \left(  \alpha^{-1}  \right)^*\omega \left(  Q,  e_1 \wedge \ldots \wedge e_n  \right)
}
{ =} { \omega \left(  \alpha^{-1} (Q) ,T_Q \left(  \alpha^{-1}  \right)(e_1) \wedge \ldots \wedge T_Q \left(  \alpha^{-1}  \right)( e_n)  \right)
}
{ } {
}
{ } {
}
{ } {
}
}
{}{}{.}
Nach Definition der metrischen Fundamentalmatrix ist

\mavergleichskettedisp
{\vergleichskette
{ g_{ij}(Q)
}
{ =} { \left\langle  T_Q \left(  \alpha^{-1}  \right) (e_i)   ,  T_Q \left(  \alpha^{-1}  \right) (e_j)   \right\rangle_{\alpha^{-1}(Q)}
}
{ } {
}
{ } {
}
{ } {
}
}
{}{}{.}
Nach
[[Euklidischer Vektorraum/Volumen eines Parallelotops/Über Skalarproduktmatrix/Fakt|Satz 7.8 (Maß- und Integrationstheorie (Osnabrück 2022-2023))]]
ist

\mavergleichskettealigndrucklinks
{\vergleichskettealigndrucklinks
{ \omega \left(  \alpha^{-1}(Q), T_Q \left(  \alpha^{-1}  \right) (e_1) \wedge \ldots \wedge T_Q \left(  \alpha^{-1}  \right) ( e_n)  \right)
}
{ =} { \left(  \det  \left(  \left\langle T_Q \left(  \alpha^{-1}  \right)  (e_i)  ,  T_Q \left(  \alpha^{-1}  \right) (e_j)   \right\rangle  \right)_{1 \leq i,j \leq n}  \right)^{1/2}
}
{ =} { \left(  \det  \left(  g_{ij}(Q)   \right)_{1 \leq i,j \leq n}  \right)^{1/2}
}
{ =} { \sqrt{g(Q)}
}
{ } {
}
}
{}{}{.}

 }

__NOEDITSECTION__

\inputaufgabeklausurloesung
{6 (1+1+3+1)}

{

Begründe, dass der $\R^2$ keine Struktur einer
\definitionsverweis {Mannigfaltigkeit mit Rand}{}{}
derart trägt, dass die angegebene Teilmenge $T$ der
\definitionsverweis {Rand}{}{}
ist.
\aufzaehlungvierabc{
\mavergleichskette
{\vergleichskette
{ T
}
{ = }{ {]0,1[}
}
{  }{
}
{    }{
}
{   }{
}
}
{}{}{,}
wobei das Intervall auf der $x$-Achse liegt.
}{
\mavergleichskette
{\vergleichskette
{ T
}
{ = }{ [0,1]
}
{  }{
}
{    }{
}
{   }{
}
}
{}{}{,}
wobei das Intervall auf der $x$-Achse liegt.
}{
\mavergleichskette
{\vergleichskette
{ T
}
{ = }{ \R
}
{  }{
}
{    }{
}
{   }{
}
}
{}{}{,}
wobei $\R$ die $x$-Achse sei.
}{
\mavergleichskette
{\vergleichskette
{ T
}
{ = }{ S^1
}
{  }{
}
{    }{
}
{   }{
}
}
{}{}{.}
}

}
{

a) Der Rand
\mathl{\partial M}{} einer Mannigfaltigkeit mit Rand ist nach
[[Mannigfaltigkeit mit Rand/Rand/Elementare Eigenschaften/Fakt|Lemma 21.3 (Differentialgeometrie (Osnabrück 2023))  (2)]]
\definitionsverweis {abgeschlossen}{}{,}
das offene Intervall als Teilmenge des $\R^2$ ist aber nicht abgeschlossen.

b) Der Rand
\mathl{\partial M}{} einer Mannigfaltigkeit mit Rand ist nach
[[Mannigfaltigkeit mit Rand/Rand ist Mannigfaltigkeit/Fakt|Satz 21.4 (Differentialgeometrie (Osnabrück 2023))]]
eine Mannigfaltigkeit ohne Rand, das abgeschlossene Intervall besitzt aber Randpunkte.

c) Es sei angenommen, dass $\R^2$ eine Mannigfaltigkeit mit dem Rand
\mathl{\R = \R \times \{0\}}{} ist. Zu
\mathl{P \in \R}{} gibt es dann eine offene Umgebung
\mathl{P\in U}{} und eine Homöomorphie
\maabbdisp {\alpha} { U } {V
} {}
mit einer offenen Menge $V \subseteq H$, wobei $H$ eine Halbebene bezeichnet. Dabei können wir $V$ durch eine kleinere offene Halbballumgebung um $P$ ersetzen. Bei einer solchen Karte werden nach
[[Mannigfaltigkeit mit Rand/Rand/Elementare Eigenschaften/Fakt|Lemma 21.3 (Differentialgeometrie (Osnabrück 2023))  (2)]]
Randpunkte auf Randpunkte abgebildet, d.h. es ist

\mavergleichskettedisp
{\vergleichskette
{\alpha ( U \cap \R  )
}
{ =} { V \cap \delta H
}
{ } {
}
{ } {
}
{ } {
}
}
{}{}{.}
Damit erhält man auch eine Homöomorphie zwischen den Komplementen, also zwischen
\mathkor {} {U \setminus U \cap \R} {und} {V \setminus V \cap \delta H} {.}
Die Halbballumgebung rechts ist aber
\definitionsverweis {wegzusammenhängend}{}{,}
wohingegen die Menge links die disjunkte Vereinigung der Schnitte mit der positiven bzw. der negativen Hälfte ist, die beide offen und auch nicht leer sind, da $U$ eine Ballumgebung von $P$ enthält. Daher ist die Menge links nicht zusammenhängend und es kann keine Homöomorphie geben.

d) wie c).

 }

__NOEDITSECTION__

\inputaufgabeklausurloesung
{4}

{

Es sei $\omega$ eine stetig differenzierbare
$n-1$-\definitionsverweis {Differentialform}{}{}
auf einer
\definitionsverweis {offenen Menge}{}{}

\mavergleichskette
{\vergleichskette
{ U
}
{ \subseteq }{ \R^n
}
{  }{
}
{    }{
}
{   }{
}
}
{}{}{}
und sei die
\definitionsverweis {äußere Ableitung}{}{}
gleich

\mavergleichskettedisp
{\vergleichskette
{d \omega
}
{ =} { fdx_1 \wedge \ldots \wedge dx_n
}
{ } {
}
{ } {
}
{ } {
}
}
{}{}{}
mit einer Funktion
\maabb {f} { U } {\R
} {.}
Zeige für

\mavergleichskette
{\vergleichskette
{ P
}
{ \in }{ U
}
{  }{
}
{    }{
}
{   }{
}
}
{}{}{}
die Gleichheit

\mavergleichskettedisp
{\vergleichskette
{ f(P)
}
{ =} { { \frac{   1  }{  \beta_n  } } \cdot \operatorname{lim}_{   \epsilon \rightarrow  0  } \, { \frac{   1  }{  \epsilon^n } } \int_{S^{n-1}(P, \epsilon)} \omega
}
{ } {
}
{ } {
}
{ } {
}
}
{}{}{,}
wobei $\beta_n$ das
[[Allgemeines Kugelvolumen/Mit Cavalieri-Prinzip/Beispiel|Volumen]]
der $n$-dimensionalen Einheitskugel bezeichnet.

}
{

Wir wenden
den [[Mannigfaltigkeit mit Rand/Satz von Stokes/Fakt|Satz von Stokes]]
auf die Kugeln
\mathl{B^n(P, \epsilon)}{} mit dem Rand
\mathl{S^{n-1}(P, \epsilon)}{}
\zusatzklammer {mit Mittelpunkt $P$ und Radius $\epsilon$} {} {} an und erhalten

\mavergleichskettealign
{\vergleichskettealign
{ \operatorname{lim}_{   \epsilon \rightarrow  0  } \, { \frac{   1 }{  \epsilon^n } }     \int_{S^{n-1}(P, \epsilon)} \omega
}
{ =} { \operatorname{lim}_{   \epsilon \rightarrow  0  } \, { \frac{   1 }{  \epsilon^n } }     \int_{ \partial B^{n}(P, \epsilon)}   \omega
}
{ =} { \operatorname{lim}_{   \epsilon \rightarrow  0  } \, { \frac{   1 }{  \epsilon^n } }     \int_{B^{n}(P, \epsilon)} d \omega
}
{ =} { \operatorname{lim}_{   \epsilon \rightarrow  0  } \, { \frac{   1 }{  \epsilon^n } }     \int_{B^{n}(P, \epsilon)} f d \lambda^n
}
{ =} {\beta_n \cdot f(P)
}
}
{}
{}{,}
wobei wir für das letzte Gleichheitszeichen die Stetigkeit von $f$ benutzt haben.

 }

__NOEDITSECTION__

\inputaufgabeklausurloesung
{12}

{

Beweise den Satz von Stokes für Mannigfaltigkeiten mit Rand.

}
{

\teilbeweis {}{}{}
{Es sei

\mathbed {U_i} {}
 {i \in I} {}
 {} {} {} {,}
eine
\definitionsverweis {offene Überdeckung}{}{}
von $M$ mit
\definitionsverweis {orientierten Karten}{}{}
und es sei

\mathbed {h_j} {}
 {j \in J} {}
 {} {} {} {,}
eine dieser
\definitionsverweis {Überdeckung untergeordnete stetig differenzierbare Partition der Eins}{}{,}
die nach
[[Mannigfaltigkeit/Abzählbare Basis/Überdeckung/Partition/Fakt|Satz 22.10 (Differentialgeometrie (Osnabrück 2023))]]
existiert. Zu jedem

\mavergleichskette
{\vergleichskette
{ P
}
{ \in }{ M
}
{  }{
}
{    }{
}
{   }{
}
}
{}{}{}
gibt es eine offene Umgebung

\mavergleichskette
{\vergleichskette
{ P
}
{ \in }{ V_P
}
{ \subseteq }{ M
}
{    }{
}
{   }{
}
}
{}{}{}
derart, dass

\mavergleichskette
{\vergleichskette
{ h_j \vert_{V_P }
}
{ = }{ 0
}
{  }{
}
{    }{
}
{   }{
}
}
{}{}{}
bis auf endlich viele Ausnahmen ist. Es sei $Y$ der Träger von $\omega$. Die Überdeckung

\mavergleichskette
{\vergleichskette
{ Y
}
{ \subseteq }{ \bigcup_{P \in M} V_P
}
{  }{
}
{    }{
}
{   }{
}
}
{}{}{}
besitzt wegen der vorausgesetzten
\definitionsverweis {Kompaktheit}{}{}
eine endliche Teilüberdeckung, sagen wir

\mavergleichskettedisp
{\vergleichskette
{ Y
}
{ \subseteq} { V_{1} \cup \ldots \cup V_{r}
}
{ =} { V
}
{ } {
}
{ } {
}
}
{}{}{.}
Daher sind überhaupt nur endlich viele der $h_j$ auf $V$ von $0$ verschieden. Wir setzen

\mavergleichskette
{\vergleichskette
{ \omega_j
}
{ = }{ h_j \omega
}
{  }{
}
{    }{
}
{   }{
}
}
{}{}{;}
diese Differentialformen sind ebenfalls stetig differenzierbar. Der Träger von $h_j$ ist eine in $M$
\definitionsverweis {abgeschlossene Teilmenge}{}{,}
die in
\mathl{U_{i(j)}}{} liegt, daher liegt der Träger von $\omega_j$ in
\mathl{Y \cap U_{i(j)}}{} und ist selbst kompakt nach
[[Kompakter Raum/Abgeschlossene Teilmenge/Kompakt/Fakt/Beweis/Aufgabe|Aufgabe 13.10 (Differentialgeometrie (Osnabrück 2023))]].
Es gilt

\mavergleichskettedisp
{\vergleichskette
{ \omega
}
{ =} { \sum_{j \in J} \omega_j
}
{ } {
}
{ } {
}
{ } {
}
}
{}{}{,}
wobei nur endlich viele dieser Differentialformen
\mathl{\omega_j}{} von $0$ verschieden sind, da

\mavergleichskette
{\vergleichskette
{ \omega_j \vert_{M \setminus Y}
}
{ = }{ 0
}
{  }{
}
{    }{
}
{   }{
}
}
{}{}{}
für alle $j$ ist und

\mavergleichskettedisp
{\vergleichskette
{ \omega_j \vert_V
}
{ =} { 0
}
{ } {
}
{ } {
}
{ } {
}
}
{}{}{}
für alle $j$ bis auf endlich viele Ausnahmen. Wegen der
[[Volumenform/Integration/Eigenschaften/Fakt|Additivität des Integrals von Differentialformen]]
und der
[[Differentialform/Mannigfaltigkeit/Äußere Ableitung/Eigenschaften/Fakt|Additivität der äußeren Ableitung]]
kann man die Aussage für die einzelnen $\omega_j$ getrennt beweisen. Wir können also annehmen, dass eine stetig differenzierbare $(n-1)$-Differentialform gegeben ist, die kompakten Träger besitzt, der ganz in einer Kartenumgebung $U$ liegt.}
{}
\teilbeweis {}{}{}
{Es liegt ein
\definitionsverweis {Diffeomorphismus}{}{}
\maabb {\alpha} { U } { V
} {}
mit

\mavergleichskette
{\vergleichskette
{ V
}
{ \subseteq }{ H
}
{  }{
}
{    }{
}
{   }{
}
}
{}{}{}
offen vor, der zugleich einen Diffeomorphismus zwischen den Rändern

\mavergleichskette
{\vergleichskette
{ \partial U
}
{ = }{ \partial M \cap U
}
{  }{
}
{    }{
}
{   }{
}
}
{}{}{}
und

\mavergleichskette
{\vergleichskette
{ \partial V
}
{ = }{ \partial H \cap V
}
{  }{
}
{    }{
}
{   }{
}
}
{}{}{}
induziert. Dabei gilt

\mavergleichskettedisp
{\vergleichskette
{
\int_{  \partial U  }   \omega
}
{ =} {
\int_{  \partial V  }  \alpha^{-1 *}\omega
}
{ } {
}
{ } {
}
{ } {
}
}
{}{}{}
und

\mavergleichskettedisp
{\vergleichskette
{
\int_{ U  }  d \omega
}
{ =} {
\int_{ V  }  \alpha^{-1 *} (  d \omega)
}
{ =} {
\int_{  V  }  d \left(  \alpha^{-1 *}\omega  \right)
}
{ } {
}
{ } {
}
}
{}{}{}
nach
[[Differentialform/Offene Menge/Äußere Ableitung/Eigenschaften/Fakt|Lemma 20.2 (Differentialgeometrie (Osnabrück 2023))  (5)]].}
{\leerzeichen{}Wir können also von einer auf

\mavergleichskette
{\vergleichskette
{ V
}
{ \subseteq }{ H
}
{  }{
}
{    }{
}
{   }{
}
}
{}{}{}
definierten stetig differenzierbaren Differentialform mit kompaktem Träger ausgehen, die wir auf ganz $H$ außerhalb des Trägers als Nullform fortsetzen können.}
\teilbeweis {}{}{}
{Wegen der Kompaktheit des Trägers gibt es einen hinreichend großen Quader

\mavergleichskette
{\vergleichskette
{ Q
}
{ \subseteq }{ H
}
{  }{
}
{    }{
}
{   }{
}
}
{}{}{,}
dessen eine Seite $S$ auf $\partial H$ liegt und der den Träger von $\omega$ nur in $S$ trifft. Auf allen anderen Seiten von $Q$ ist $\omega$
\zusatzklammer {und damit auch $d\omega$} {} {}
die Nullform. Daher gilt einerseits

\mavergleichskettedisp
{\vergleichskette
{
\int_{ H  }  d\omega
}
{ =} {
\int_{ Q  }  d\omega
}
{ } {
}
{ } {
}
{ } {
}
}
{}{}{}
und andererseits

\mavergleichskettedisp
{\vergleichskette
{
\int_{  \partial H  }   \omega
}
{ =} {
\int_{  S  }   \omega
}
{ =} {
\int_{  \partial Q  }   \omega
}
{ } {
}
{ } {
}
}
{}{}{.}
Somit folgt die Aussage
aus der [[Satz von Stokes/Quaderversion/Fakt|Quaderversion des Satzes von Stokes]].}
{}

 }
 [[Kategorie:Latexseite]]
